\sscp{\tsc{Patakai-Manusela}}{10-03-03}
\citet{Stresemann1927} did not define his Sub-Seran subgroup
other than as Proto-Ambon languages that did not fit with Sub-Ambon or Sub-Buru.
Similarly, \citet{Collins1983} did not discuss the innovations
defining the \tsc{Patakai-Manusela} branch of his Nunusaku (roughly our \tsc{Ambon-Seram}).
This gap in our understanding encourages us to look at these languages more closely.
From \citet{Collins1983}, \tsc{Patakai-Manusela} includes the
following languages: the Manusela cluster,
North Nuaulu, and South Nuaulu (also known as “Patakai” or “Fatakai”).

\begin{table}
	\centering\tcp{ Lexical similarity among varieties of Manusela}{10-10}
	\st{6pt}\begin{tabular}{llllllllll}\lsptoprule
\mc{4}{l}{\it{Alakmat} {\cgr}}									 		 		&{\cgr}		&{\cgr}		&{\cgr}		&{\cgr}		&{\cgr}		& Huaulu [hud]\\
72{\cgr}	&	\mc{4}{l}{\it{Kanikeh}}								 		 		&		&		&		&		& Kanikeh\\	\hhline{~~--------}
66{\cgr}	&	67	&	\mc{4}{|l}{\it{Wanesa} {\cgr}}								 		 	&{\cgr}		&{\cgr}		&{\cgr}		&	\mc{1}{l|}{}\\	
72{\cgr}	&	74	&	\mc{1}{|l}{80{\cgr}}	&	\mc{4}{l}{\it{Piliana}}						 		 		&		&		& \mc{1}{l|}{South Manusela}\\
69{\cgr}	&	72	&	\mc{1}{|l}{76{\cgr}}	&	79	&	\mc{4}{l}{\it{Yamahena}{\cgr}}					 		 		&{\cgr}		&	\mc{1}{l|}{}\\	\hhline{~~--------}
71{\cgr}	&	74	&	67{\cgr}	&	73	&	72{\cgr}	&	\mc{3}{l}{\it{Air Besar}}			 		&		& Hatuolo\\	\hhline{~~~~~~----}
64{\cgr}	&	66	&	65{\cgr}	&	70	&	72{\cgr}	&	75	&	\mc{3}{|l}{\it{Maneo Rendah}{\cgr}}	 			&\mc{1}{l|}{}\\
66{\cgr}	&	67	&	66{\cgr}	&	72	&	73{\cgr}	&	73	&	\mc{1}{|l}{86{\cgr}}	&	\mc{2}{l}{\it{Maneoratu}}	&\mc{1}{l|}{\multirow{-2}{*}{Maneo}}\\	\hhline{~~~~~~----}
		\lspbottomrule
	\end{tabular}
\end{table}

Evidence in \citet[117--121]{LoskiLoski1989} suggests that “Manusela” (a.k.a.
Sou Upaa) is best thought of as a chain or cluster of related
languages and dialects, rather than a single language, with the different
varieties surveyed averaging 71{\%} lexical similarity.
Of these varieties, only Huaulu (originally spoken in the interior)
has been assigned a separate ISO 639-3 code.
This variety is also referred to by other Manusela speakers as
the \it{bahasa asli} ‘native/original language’.
The percentages of lexical similarities calculated by \citet[119]{LoskiLoski1989}
are given in \trf{10-10} along with the villages surveyed in
italics and the names of the varieties
identified by \citet{LoskiLoski1989} in normal typeface.

Data throughout this section is taken from the Nusaweleh variety,
documented by \citet{MacDonald2015} which is spoken in the central
northern area marked “Manusela cluster” in \frf{07-03}.
Data marked “(K)” is the Kowa variety of Manusela from the Holle list \citep[185--191]{Stokhof1981a}.
South Nuaulu is currently spoken in a few villages in south central Seram.
It is related to but distinct from North Nuaulu,
spoken near the north coast of central Seram.

\tsc{Patakai-Manusela} shares the merger of *n/*ñ/*ŋ  > \it{n},
weakening of *p, and *mb/*mp > **mp which are the changes which
differentiate \tsc{Ambon-Seram} from neighbouring subgroups.
They also share merger of *R/*l/*j > **l which occurs in all
\tsc{Ambon-Seram} languages except Boano (\srf{10-03-01}).
However, they have distinct *d/*z > **h that sets them apart
as a separate branch within \tsc{Ambon-Seram}.
Examples of *ŋ > \it{n} in are given in \qf{10-47},
and examples of *ñ > \it{n} are given in \qf{10-48}.

\noindent{\wg\st{2.5pt}\begin{tabularx}
{\linewidth}{>{\hspace{-\tabcolsep}}l<{\hspace{-\tabcolsep}}
>{\hspace{-\tabcolsep}\enspace}llll
>{\raggedright\arraybackslash}X}
\lsptoprule
%\tbx{10-47}	&	PMP	&	Manusela	&	S. Nuaulu	&	N. Nuaulu	&	local gl.	\\	\midrule
%\endfirsthead
	&	PMP	&	Manusela	&	S. Nuaulu	&	N. Nuaulu	&	local gl.	\\	\midrule
%\endhead
\tbx{10-47}	&	*\tbr{ŋ}ajan	&	\tbr{l}ala-ni	&	\tbr{n}ana-i	&	\tbr{n}ana-e	&	name	\\
	&	*\tbr{ŋ}ilu/*\tbr{n}ilu	&	ai aka maka\tbr{n}inu-a	&	ma\tbr{n}inu-e	&	ma\tbr{n}inu-e	&	Man. `tamarind', S. Nuaulu `sour'	\\
	&	*\tbr{ŋ}is\tbr{ŋ}is	&	\tbr{n}esi-ni	&	\tbr{n}esi-ni	&	\tbr{n}esi-ri	&	teeth	\\
	&	**\tbr{ŋ}əli{\footnotemark} 	&		&	\tbr{n}eni-ni	&		&	canine tooth	\\
	&	*tali\tbr{ŋ}a	&	ti\tbr{n}a-ni-a	&	ti\tbr{n}a-i	&	tiri\tbr{n}a	&	ear	\\
	&	*təri\tbr{ŋ}	&	tehi	&	ten\tbr{n}-e	&		&	giant bamboo	\\
	&	*ha\tbr{ŋ}in	&	a\tbr{n}in-o	&	\tcg{(ihut-e)}	&	\tcg{(ifut-e)}	&	wind	\\
	&	*na\tbr{ŋ}uy	&	na\tbr{n}u (K)	&	akina\tbr{n}u	&	yakina\tbr{n}u	&	swim	\\
	&	*sa\tbr{ŋ}a	&	sa\tbr{n}ahata	&	sa\tbr{n}a-e	&		&	branch	\\
	&	*sa\tbr{ŋ}əlaR	&	(korenia)	&	asi\tbr{n}ina	&		&	fry	\\
	&	*ta\tbr{ŋ}is	&	puta\tbr{n}i	&	ra\tbr{n}i	&		&	cry	\\
	&	*ba\tbr{ŋ}un	&	pa\tbr{n}u	&	ha\tbr{n}u	&		&	get/wake up	\\
	&	*pi\tbr{ŋ}us	&	‹i\tbr{n}oë› (K)	&	\tcg{(hunn-e)}	&	\tcg{(nofonafuaya)}	&	snot	\\
	&	*la\tbr{ŋ}it	&	la\tbr{l}i-a	&	na\tbr{n}t-e	&	na\tbr{n}at-e	&	sky	\\	\midrule
\tbx{10-48}	&	*ma-qa\tbr{ñ}ud	&		&	ma\tbr{n}u	&		&	float, drift	\\
	&	*pə\tbr{ñ}u	&	pe\tbr{n}u-a (K)	&	e\tbr{n}u	&		&	turtle	\\
	&	*mi\tbr{ñ}ak	&		&	mi\tbr{n}a-i	&		&	meat, fat, oil	\\
\lspbottomrule
\end{tabularx}\mbox{}\\}
\footnotetext{For *ŋəli ‘canine tooth’ compare Kamarian \it{neri}
	‘tusks of pig’ \citep[315]{vanEkris1864},
	Hitu \it{neli} ‘tusk (of a pig)’ \citep[105]{vanHoevell1877}.
	For initial *ŋ compare cognates in Timor: Helong \it{ŋilin},
	Meto \it{niis} \it{koni-f} ‘canine tooth’,
	Proto-Rote-Meto *ŋoli \citep[297]{Edwards2021}.}

The reflexes of *ŋajan > \it{lala-ni} ‘name’
and *laŋit > \it{lali-a} in Manusela appear to show progressive
and regressive assimilation respectively of the nasal to the lateral.
While such assimilation is frequent among \tsc{Ambon-Seram} languages
(see the discussion of Boano in \srf{10-03-01}
and \tsc{Atamanu} in \srf{10-03-02}) there are,
unfortunately no comparable reflexes of words with *nVl or *lVn
in available Manusela data to test whether this is a regular pattern.
Note also that such putative assimilation does \it{not} occur
in *ŋilu/*nilu > \it{ai aka makaninu} ‘tamarind’.
Instead, this word appears to show assimilation of the lateral to the nasal.

Examples of *p > \it{h} > {\0} are given in \qf{10-49}.
\\

\noindent{\wg\st{6pt}\begin{tabularx}
{\linewidth}{>{\hspace{-\tabcolsep}}l<{\hspace{-\tabcolsep}}
>{\hspace{-\tabcolsep}\enspace}llll
>{\raggedright\arraybackslash}X}
\lsptoprule
%\tbx{10-49}	&	PMP	&	Manusela	&	S. Nuaulu	&	N. Nuaulu	&	local gl.	\\	\midrule
%\endfirsthead
\tbx{10-49}	&	PMP	&	Manusela	&	S. Nuaulu	&	N. Nuaulu	&	local gl.	\\	\midrule
%\endhead
	&	*\tbr{p}ajay	&	\tbr{h}ala	&	ana	&		&	rice (Kowa Man. \it{fala})	\\
	&	*\tbr{p}ija	&	\tbr{h}ila	&	ina	&	ina	&	how much?	\\
	&	*\tbr{p}itu	&	\tbr{h}itu	&	itu	&	itu	&	seven	\\
	&	*\tbr{p}iliq	&	li	&	nini	&		&	choose	\\
	&	*\tbr{p}anah	&		&	ana	&	\tcg{(iketa)}	&	shoot	\\
	&	*ə\tbr{p}at	&	\tbr{h}ate	&	ate	&	ate	&	four	\\
	&	**qali\tbr{p}an	&	ni{\tl}ni\tbr{h}al-a	&	ni{\tl}nian-e	&	ni{\tl}nian-e	&	centipede (C metathesis in Manusela)	\\
	&	*qa\tbr{p}əju	&	wolu-ni	&	\tcg{(moni-ni)}	&	\tcg{(moni-n)}	&	bile	\\
	&	*ma-hi\tbr{p}i	&	mi-nia	&	niin-e	&		&	dream	\\
	&	*də\tbr{p}a	&		&	noa	&		&	fathom (irr. *d > \it{n})	\\
\lspbottomrule
\end{tabularx}\mbox{}\\}

Examples of *ma-b/*ma-p/*mb/*mp > \it{p} are given in \qf{10-50}.
South Nuaulu \it{mpia} ‘raw sago’ probably attests usual *R >
**n (see \qf{10-54} below) with subsequent assimilation of **n
to the place of the following plosive.
That is, something like *Rumbia > **Rupia > **nupia > **npia
> \it{mpia}.

\noindent{\wg\st{3pt}\begin{tabularx}
{\linewidth}{>{\hspace{-\tabcolsep}}l<{\hspace{-\tabcolsep}}
>{\hspace{-\tabcolsep}\enspace}llll
>{\raggedright\arraybackslash}X}
\lsptoprule
%\tbx{10-50}	&	PMP	&	Manusela	&	S. Nuaulu	&	N. Nuaulu	&	local gl.	\\	\midrule
%\endfirsthead
\tbx{10-50}	&	PMP	&	Manusela	&	S. Nuaulu	&	N. Nuaulu	&	local gl.	\\	\midrule
%\endhead
	&	*\tbr{ma-p}utiq	&	\tbr{p}uti-a	&	\tbr{p}uti-e	&	\tbr{p}uti-e	&	white	\\
	&	*\tbr{ma-p}ənuh	&	\tbr{p}onu-i	&	\tcg{(tau-e)}	&	\tcg{(tau-e)}	&	full	\\
	&	**\tbr{ma-p}utu{\footnotemark} 	&		&	\tbr{p}utu	&		&	hot (blood)	\\
	&	*ka\tbr{mb}u	&	a\tbr{p}u	&	a\tbr{p}u-i	&	a\tbr{p}u-ri	&	stomach, belly	\\
	&	*Ru\tbr{mb}ia	&	i\tbr{p}ia	&	\tbr{mp}ia-e, \tbr{p}ia-e	&	\tcg{(fatanetinai)}	&	Man. `sago' S. Nuaulu `raw sago'	\\
	&	**ka\tbr{mb}əR{\footnotemark} 	&		&	a\tbr{p}ana	&	rea\tbr{p}ana	&	saliva	\\
	&	*u\tbr{mp}u	&	u\tbr{p}u (K),	&	u\tbr{p}u	&	u\tbr{p}u	&	ancestor, descendant \\ \hhline{~}
	&							&	u\tbr{h}u-na	&							&				&	(\it{uhu-na} < *upu)	\\
\lspbottomrule
\end{tabularx}\mbox{}\\}
\footnotetext[10]{For **(ma-)putu compare Buru \it{poto} ‘hot,
	burn’, Kayeli \it{mpoto} ‘fever’, Proto-Rote-Meto *mbutu ‘hot’
\citep[273]{Edwards2021}.}
\footnotetext{For **kambəR compare Kamarian \it{aper} \citep[76]{vanEkris1864} among others.
	For medial **mb compare Proto-Rote-Meto *kambe ‘saliva,
	spittle’ and subentries in \citet[192]{Edwards2021}.}

\tsc{Patakai-Manusela} has *d/*z > **h.
This is retained as \it{h} in Manusela
and is lost in Nuaulu.
Examples are given in \qf{10-51}
and \qf{10-52}.
\\

\noindent{\wg\st{3pt}\begin{tabularx}
{\linewidth}{>{\hspace{-\tabcolsep}}l<{\hspace{-\tabcolsep}}
>{\hspace{-\tabcolsep}\enspace}llll
>{\raggedright\arraybackslash}X}
\lsptoprule
%\tbx{10-51}	&	PMP	&	Manusela	&	S. Nuaulu	&	N. Nuaulu	&	local gl.	\\	\midrule
%\endfirsthead
	&	PMP	&	Manusela	&	S. Nuaulu	&	N. Nuaulu	&	local gl.	\\	\midrule
%\endhead
\tbx{10-51}	&	*\tbr{d}uha	&	\tbr{h}ua	&	ua	&	ua	&	two	\\
	&	*\tbr{d}uRi	&	\tbr{h}ulikata	&	uni-e	&	uni-e	&	bone	\\
	&	*\tbr{d}aləm	&	ria \tbr{h}ali	&	ano-i	&	ano-e	&	inside	\\
	&	*\tbr{d}akəp{\footnotemark} 	&	\tbr{h}aku	&	\tcg{(nete, ai)}	&		&	catch, grab, hold, take	\\
	&	*kə\tbr{d}əŋ	&	ko\tbr{h}o	&	oo	&	yoo	&	stand	\\
	&	*ana\tbr{d}uq	&	\tcg{(meleka)}	&	nau-e	&	nau-e	&	long	\\
	&	*ma-u\tbr{d}əhi	&	\tcg{(kalu-nia)}	&	mui-e	&	mui-e	&	behind, last	\\
	&	*tuhu\tbr{d}	&	tu\tbr{h}e-ni	&	tune-i	&	tune-ri	&	knee	\\	\midrule
\tbx{10-52}	&	*\tbr{z}alan	&	\tbr{h}ala	&	\tcg{(arena)}	&	\tcg{(arena)}	&	path, way	\\
	&	*qu\tbr{z}an	&	\tcg{(roa)}	&	uan-e	&	uan-e	&	rain	\\
	&	*haRə\tbr{z}an	&	o\tbr{h}a	&	oan-e	&	\tcg{(autanat-e)}	&	ladder	\\
	&	*tu\tbr{z}uq	&	katu\tbr{h}u	&	au-tuu	&		&	Man. `clue, hint' S.N. `point'	\\
	&	*qa\tbr{z}ay	&	a\tbr{h}a-nia	&	a\tbr{n}a-i	&	a\tbr{n}a-ri nofua	&	chin, jaw, cheek (irr. *z > \it{n} in Nuaulu)	\\
\lspbottomrule
\end{tabularx}\mbox{}\\}
\footnotetext{South Nuaulu \it{ai} ‘take’ is possibly a reflex of
*dakəp, but this requires irregular *ə > \it{i}.}

Four exceptions have been identified in South Nuaulu.
Three show *d/*z > \it{n} *qazay > \it{ana-i} ‘chin’,
and *kudəŋ > \it{unen-e} ‘pot with a narrow neck’,
and *dəpa > \it{noa} ‘fathom’.
The other possible exception is *tuqəd > \it{tuete} ‘tree stump’
which either shows *d > \it{t} /{\gap}{\#},
or has the nominal suffix \it{-te}.
One possible exception has been identified in Manusela *zəlay
(PWMP) > \it{lela-ni} ‘tongue’.

Like many \tsc{Ambon-Seram} languages,
Manusela and Nuaulu share *R/*l/*j > **l.
This is retained as \it{l} in Manusela
and shifts to \it{n} in Nuaulu.
Examples are given in \qf{10-53}--\qf{10-55}.
\\

\noindent{\wg\st{2.25pt}\begin{tabularx}
{\linewidth}{>{\hspace{-\tabcolsep}}l<{\hspace{-\tabcolsep}}
>{\hspace{-\tabcolsep}\enspace}l<{\hspace{-\tabcolsep}\,}>{\hspace{-\tabcolsep}\,}llll
>{\raggedright\arraybackslash}X}
\lsptoprule
%\tbx{10-53}	&		&	P\tcg{(CE)}MP	&	Manusela	&	S. Nuaulu	&	N. Nuaulu	&	local gl.	\\	\midrule
%\endfirsthead
&		&	PMP	&	Manusela	&	S. Nuaulu	&	N. Nuaulu	&	local gl.	\\	\midrule
%\endhead
\tbx{10-53}		&	*l	&	*\tbr{l}ima	&	\tbr{l}ima	&	\tbr{n}ima	&	\tbr{n}ima	&	five	\\
	&		&	*\tbr{l}iqəR	&	\tbr{l}io-na	&	\tbr{n}ion-i, nio-i	&	\tbr{n}ioka	&	voice, PMP `neck'	\\
	&		&	*\tbr{l}ahud	&	\tcg{(meleke)}	&	hai\tbr{n}au	&	fai\tbr{n}au	&	far (PMP `sea')	\\
	&		&	*ba\tbr{l}u	&	ha\tbr{l}u-a	&	uha\tbr{n}u	&		&	widow(er)	\\
	&		&	*bu\tbr{l}u	&	hu\tbr{l}u-ni	&	hu\tbr{n}u-e	&	hu\tbr{n}u-a	&	body hair	\\
	&		&	*tə\tbr{l}u	&	to\tbr{l}u	&	to\tbr{n}u	&	to\tbr{n}u	&	three	\\
	&		&	*qu\tbr{l}əj	&	u\tbr{l}eh	&	u\tbr{n}e	&	u\tbr{n}e putut-e	&	maggot N. Nuaulu `caterpillar'	\\
	&		&	*wa\tbr{l}u	&	wa\tbr{l}u	&	wa\tbr{n}u	&	wa\tbr{n}u	&	eight	\\
	&		&	**ma\tbr{l}ip	&	muma\tbr{l}i	&	au-ma\tbr{n}i	&	yau-ma\tbr{n}	&	laugh	\\
	&		&	*bu\tbr{l}an	&	\tcg{(umala)}	&	hu\tbr{n}an-e	&	hu\tbr{n}an	&	moon	\\
	&		&	*qatə\tbr{l}uR	&	to\tbr{l}u-a	&	toun-e	&	\tcg{(etut-e)}	&	egg (irr. *l > {\0} in S.N.)	\\
	&		&	**ma\tbr{l}abaw	&	ma\tbr{i}aha	&	m\tbr{n}aha	&	m\tbr{n}afa	&	mouse, rat (irr. *l > \it{i} in Man.)	\\	\midrule
\tbx{10-54}	&	*R	&	*\tbr{R}umaq	&	\tbr{l}uma	&	\tbr{n}uma	&	\tbr{n}uma	&	house	\\
	&		&	*\tbr{R}usuk	&	\tbr{n}usuala	&	\tbr{n}usu-i	&	\tcg{(nureana-e)}	&	ribs (irr. *R > \it{n} in Man.)	\\
	&		&	*\tbr{R}aməs	&	\tcg{(poso)}	&	\tbr{n}ama	&		&	squeeze (irr. *ə > \it{a})	\\
	&		&	*\tbr{R}ibu	&	\tbr{l}ihuni	&	\tbr{n}ihun	&	\tbr{n}ifun	&	thousand	\\
	&		&	*\tbr{R}amut	&	\tbr{l}amu-nia	&	\tbr{n}amt-e	&	\tbr{n}amt-e	&	Man. `fibre', S. Nuaulu `root'	\\
	&		&	*qaba\tbr{R}a	&	ha\tbr{l}a (K)	&	ha\tbr{n}a-i	&	fa\tbr{n}a-ri 	&	arm PMP shoulder	\\
	&		&	*baqə\tbr{R}u	&	ho\tbr{l}u-a	&	ho\tbr{n}u-e	&	yaho\tbr{n}u	&	new	\\
	&		&	*u\tbr{R}at	&	\tcg{(loloinya)}	&	u\tbr{n}at-e	&	u\tbr{n}at-a	&	vein	\\
\lspbottomrule
\end{tabularx}\mbox{}\\}

\noindent{\wg\st{2.25pt}\begin{tabularx}
{\linewidth}{>{\hspace{-\tabcolsep}}l<{\hspace{-\tabcolsep}}
>{\hspace{-\tabcolsep}\enspace}l<{\hspace{-\tabcolsep}\,}>{\hspace{-\tabcolsep}\,}llll
>{\raggedright\arraybackslash}X}
\lsptoprule
%\tbx{10-53}	&		&	P\tcg{(CE)}MP	&	Manusela	&	S. Nuaulu	&	N. Nuaulu	&	local gl.	\\	\midrule
%\endfirsthead
&		&	PMP	&	Manusela	&	S. Nuaulu	&	N. Nuaulu	&	local gl.	\\	\midrule
%\endhead
\tbx{10-55}	&	*j	&	*pa\tbr{j}ay	&	ha\tbr{l}a	&	a\tbr{n}a	&		&	rice	\\
	&		&	*pi\tbr{j}a	&	hi\tbr{l}a	&	i\tbr{n}a	&	i\tbr{n}a	&	how much?	\\
	&		&	*ŋa\tbr{j}an	&	la\tbr{l}a-ni	&	na\tbr{n}a-i	&	na\tbr{n}a-e	&	name	\\
	&		&	*ma\tbr{j}a	&	muma\tbr{l}a	&	mama\tbr{n}a-e	&	remama\tbr{n}a-e	&	dry	\\
	&		&	*qapə\tbr{j}u	&	wo\tbr{l}u-ni	&	\tcg{(moni-ni)}	&	\tcg{(moni-n)}	&	bile	\\
	&		&	*hua\tbr{j}i	&		&	wa\tbr{n}i-ni	&	wa\tbr{n}i-n	&	S. Nuaulu `same sex ySib', N. Nuaulu `grandchild' (gloss error?)	\\
\lspbottomrule
\end{tabularx}\mbox{}\\}

Manusela tends to have word-final *R > {\0} while South Nuaulu
tends to retain word-final *R as \it{n}.
This is consistent with the tendency for South Nuaulu to preserve word final consonants.
Examples in addition to *qatəluR
and *liqəR in \qf{10-53} above  include *wahiyəR > Manusela \it{wae-a}
South Nuaulu, North Nuaulu \it{waen-e}, *niuR/**niwəR > Manusela \it{nue-a},
South Nuaulu \it{nion-e} and *mantəR > South Nuaulu \it{maran-e} ‘cuscus’.
Both languages usually have word final *j > {\0}.
Data for word-final *l is too sparse to identify regular patterns.

A few words have irregular *l > {\0}.
Such words include *kalaw > Manusela \it{kaa} ‘Blyth’s hornbill
(Rhyticeros plicatus)’, *qalima > Manusela \it{ima-ni} ‘hand,
wrist’, and *qatəluR > South Nuaulu \it{toun-e} ‘egg’ (Manusela \it{tolu-a}).
Consistent with other \tsc{Ambon-Seram} languages (see \trf{07-05},
\srf{07-03-02}) *qalejaw  has *j > {\0} in Manusela \it{lea} ‘sun’.
(South Nuaulu has non-cognate \it{rani-e} ‘sun’).

The normal reflex of schwa in Manusela
and Nuaulu is \it{o},
as shown in \qf{10-56}.
Some words show *ə > \it{e} or *ə > \it{i},
as shown in \qf{10-57}.
In most cases this appears to be conditioned by a following high
vowel or palatal consonant,
including final *əj > \it{e}.
Many of these words also show *ə > e in other \tsc{Ambon-Seram}
languages which otherwise have *ə > \it{o}.
Thus, for instance, reflexes of *ma-qitəm ‘black’ has final *ə
> \it{e} in all \tsc{Ambon-Seram} languages for which data is
available, including those which usually have other reflexes of *ə in final syllables.
Indeed, *ə > \it{e} in this word can be identified as a lexically
specific innovation at the level of Proto-Ambon-Seram.

\noindent{\wg\st{2.5pt}\begin{tabularx}
{\linewidth}{>{\hspace{-\tabcolsep}}l<{\hspace{-\tabcolsep}}
>{\hspace{-\tabcolsep}\enspace}llll
>{\raggedright\arraybackslash}X}
\lsptoprule
%\tbx{10-56}	&	PMP	&	Manusela	&	S. Nuaulu	&	N. Nuaulu	&	local gl.	\\	\midrule
%\endfirsthead
	&	PMP	&	Manusela	&	S. Nuaulu	&	N. Nuaulu	&	local gl.	\\	\midrule
%\endhead
\tbx{10-56}	&	*t\tbr{ə}buh	&	t\tbr{o}hu-a	&	t\tbr{o}hu	&	t\tbr{o}hu	&	sugar cane	\\
	&	*b\tbr{ə}taw	&	h\tbr{o}ta (K)	&	h\tbr{o}ta	&	f\tbr{o}ta-i	&	opposite sex sibling	\\
	&	*b\tbr{ə}t\tbr{ə}ŋ	&		&	h\tbr{o}t\tbr{o}n-e	&		&	kind of grain PMP `foxtail millet'	\\
	&	*hut\tbr{ə}k	&		&	ut\tbr{o}-e	&		&	marrow, forehead, brain	\\
	&	*baq\tbr{ə}Ru	&	h\tbr{o}lu-a	&	h\tbr{o}nu-e	&	yah\tbr{o}nu	&	new	\\
	&	*liq\tbr{ə}R	&	li\tbr{o}-na	&	ni\tbr{o}n-i, ni\tbr{o}-i	&	ni\tbr{o}ka	&	voice, PMP neck	\\
	&	*t\tbr{ə}lu	&	t\tbr{o}lu	&	t\tbr{o}nu	&	t\tbr{o}nu	&	three	\\
	&	*haR\tbr{ə}zan	&	\tbr{o}ha	&	\tbr{o}an-e	&	\tcg{(autanat-e)}	&	ladder	\\
	&	*qap\tbr{ə}ju	&	w\tbr{o}lu-ni	&	\tcg{(moni-ni)}	&	\tcg{(moni-n)}	&	bile	\\
	&	*qat\tbr{ə}luR	&	t\tbr{o}lu-a	&	t\tbr{o}un-e	&	\tcg{(etut-e)}	&	egg (S.N. irr. *l > {\0})	\\
	&	*k\tbr{ə}d\tbr{ə}ŋ	&	k\tbr{o}h\tbr{o}	&	\tbr{oo}	&	y\tbr{oo}	&	stand	\\
	&	*d\tbr{ə}pa	&		&	n\tbr{o}a	&		&	fathom (irr. *d > \it{n})	\\	\midrule
\tbx{10-57}	&	*p\tbr{ə}ñu	&	p\tbr{e}nu-a (K)	&	\tbr{e}nu	&		&	turtle	\\
	&	*t\tbr{ə}riŋ	&	t\tbr{e}hi	&	t\tbr{e}nn-e	&		&	giant bamboo	\\
	&	*b\tbr{ə}li	&		&	ah\tbr{e}ni	&	yaf\tbr{e}nike	&	sell	\\
	&	*ul\tbr{ə}j	&	ul\tbr{e}h	&	un\tbr{e}	&	u\tbr{n}e putut-e	&	maggot, N. Nuaulu `caterpillar'	\\
	&	*waR\tbr{ə}j	&	w\tbr{e}li-nia	&	a-wan\tbr{e}	&	auwan\tbr{e}	&	vine, rope	\\
	&	*tuq\tbr{ə}d	&		&	tu\tbr{e}te	&		&	tree stump (*əd{\#} > \it{e}?, or nominal suffix \it{-te})	\\
	&	*qat\tbr{ə}p	&	wat\tbr{e}-a	&	ainat\tbr{a}-e	&		&	thatch, roof (S. Nuaulu morphology unclear, possibly non-cognate)	\\
	&	*ma-qit\tbr{ə}m	&	met\tbr{e}-a, meta	&	met\tbr{e}n-e	&	met\tbr{e}n-e	&	black	\\
\lspbottomrule
\end{tabularx}\mbox{}\\}

To summarise,
Manusela and Nuaulu can be placed within \tsc{Ambon-Seram} on
the basis of shared *ŋ/*ñ/*n > \it{n},
*p > **h, and *ma-p/*mb/*mp > \it{p}.
In addition they share *R/*l/*j > **l with most other \tsc{Ambon-Seram} languages.
However, they have a distinct reflex of *d/*z > **h,
and on this basis we place \tsc{Patakai-Manusela} as a primary
branch within \tsc{Nuclear Ambon-Seram}.
